% ============================================================
%  TEMA 3 — APARTADO 5: PROCESOS SUSTRACTIVOS
%                        (Demanda < Oferta: exceso de plantilla)
%
%  RECORDATORIO AL EQUIPO:
%  Los procesos sustractivos son aquellos mediante los cuales la organización
%  reduce su plantilla cuando la oferta de RRHH supera la demanda.
%  Se clasifican en:
%
%  A) SUSPENSIÓN LABORAL TEMPORAL → ERTES
%     - Expediente de Regulación Temporal de Empleo.
%     - El contrato se SUSPENDE temporalmente (no se extingue).
%     - El trabajador cobra prestación del SEPE + posible complemento empresa.
%     - Causas: económicas, técnicas, organizativas, productivas, o FUERZA MAYOR.
%     - Muy usado durante la pandemia COVID-19.
%     - Ventaja: mantiene la plantilla formada cuando mejore la situación.
%     - Inconveniente: coste económico y de imagen; desmotivación.
%
%  B) RUPTURA LABORAL VOLUNTARIA
%     1. Dimisión (renuncia voluntaria del trabajador).
%        - El empleado abandona sin coste para la empresa (salvo sustitución).
%        - Puede ser signo de insatisfacción o de mejores oportunidades externas.
%     2. Jubilación (ordinaria, anticipada o parcial).
%        - Jubilación anticipada: la empresa incentiva la salida antes de la
%          edad legal, normalmente con complementos económicos.
%        - Es la forma más suave de reducir plantilla.
%
%  C) RUPTURA LABORAL INVOLUNTARIA
%     1. Despido individual (por causas objetivas o disciplinarias).
%     2. ERE — Expediente de Regulación de Empleo (despidos colectivos).
%        - Requiere periodo de consultas con representantes sindicales.
%        - Indemnizaciones + posibles medidas sociales.
%     3. Downsizing: reducción sistemática del tamaño organizativo.
%        - Puede acompañarse de reestructuración (rightsizing).
%     4. No renovación de contratos temporales.
%
%  D) RECOLOCACIÓN — OUTPLACEMENT
%     - Servicio de apoyo a los empleados que salen para que encuentren
%       nuevo empleo (orientación profesional, CV, entrevistas, networking).
%     - Puede ser gestionado internamente o a través de consultoras.
%     - Mejora la imagen de la empresa y reduce el trauma de la salida.
%
%  PARA ESTA SECCIÓN INDICAD:
%  - ¿Ha utilizado CaixaBank cada uno de estos mecanismos?
%  - Cuándo, bajo qué circunstancias y con qué resultados.
%  - Ventajas e inconvenientes de cada uno en el caso de CaixaBank.
%  - ¿Está la empresa satisfecha con cómo los ha gestionado?
% ============================================================

\section{Procesos sustractivos: gestión del exceso de plantilla}

Cuando la oferta de recursos humanos supera la demanda, CaixaBank debe
recurrir a mecanismos de reducción o reajuste de plantilla. A continuación
se analizan los principales procesos sustractivos aplicados o aplicables.

\subsection{Suspensión laboral temporal: los ERTE}

% TODO: ¿Ha recurrido CaixaBank a ERTE? Principalmente durante la pandemia
% COVID-19 (2020-2021). Si es así, ¿a cuántos empleados afectó? ¿Qué duración
% tuvo? ¿Qué complementos ofreció la empresa sobre la prestación del SEPE?
% ¿Fue eficaz para mantener el empleo y la operatividad del banco?

Un ERTE (Expediente de Regulación Temporal de Empleo) implica la suspensión
temporal del contrato de trabajo sin que se produzca su extinción. El
trabajador cobra una prestación por desempleo del SEPE, que la empresa puede
complementar.

\subsubsection{Aplicación en CaixaBank}

[Descripción a completar]

\textbf{Ventajas del ERTE para CaixaBank:}
\begin{itemize}
  \item [...]
\end{itemize}

\textbf{Inconvenientes / limitaciones:}
\begin{itemize}
  \item [...]
\end{itemize}

\subsection{Ruptura laboral voluntaria}

\subsubsection{Dimisión y rotación voluntaria}

% TODO: ¿Cuál es la tasa de dimisiones voluntarias en CaixaBank?
% ¿Realizan entrevistas de salida para conocer las causas?
% ¿Qué factores llevan a los empleados a marcharse voluntariamente
% (mejores salarios en la competencia, agotamiento, cambio de sector...)?

[Descripción a completar]

\subsubsection{Jubilación (ordinaria y anticipada)}

% TODO: CaixaBank tiene una plantilla relativamente envejecida en ciertos
% segmentos (herencia de Bankia y de la red de oficinas tradicionales).
% ¿Han implementado planes de jubilación anticipada?
% ¿En el ERE de 2021 se incluyeron medidas de jubilación anticipada?
% ¿Con qué condiciones (edad mínima, complemento sobre la pensión)?

[Descripción a completar]

\textbf{Ventajas de los planes de jubilación anticipada:}
\begin{itemize}
  \item Reducción no traumática de la plantilla sin despidos.
  \item Renovación generacional y entrada de perfiles más digitales.
  \item Menor conflicto sindical que los despidos colectivos.
\end{itemize}

\textbf{Inconvenientes:}
\begin{itemize}
  \item Pérdida de conocimiento experto y memoria institucional.
  \item Coste económico elevado para la empresa (complementos).
  \item Riesgo de que se jubilen empleados clave que no deberían marcharse.
\end{itemize}

\subsection{Ruptura laboral involuntaria: ERE y despidos}

\subsubsection{ERE 2021 — Fusión CaixaBank-Bankia}

% TODO: El ERE de 2021 fue el más grande de la historia bancaria española.
% Detalles clave a incluir:
%   - Número de afectados (aproximadamente 6.452 empleados según acuerdo sindical).
%   - Condiciones: indemnizaciones, bajas incentivadas, jubilaciones anticipadas.
%   - Sindicatos involucrados: UGT, CCOO, FINE.
%   - Reacción de los empleados y sindicatos.
%   - ¿Fue el ajuste suficiente o hay nuevos ajustes previstos?
%   - ¿Cómo afectó a la moral de la plantilla que se quedó?

En 2021, tras la fusión con Bankia, CaixaBank llevó a cabo el mayor
Expediente de Regulación de Empleo (ERE) de la banca española:

\begin{table}[H]
  \centering
  \small
  \caption{Datos del ERE CaixaBank 2021}
  \begin{tabularx}{\textwidth}{lX}
    \toprule
    \textbf{Parámetro}           & \textbf{Dato} \\
    \midrule
    Empleados afectados           & $\sim$6.452 (según acuerdo negociado) \\
    Cierre de oficinas            & $\sim$1.500 oficinas \\
    Tipo de salidas               & Bajas incentivadas, jubilaciones anticipadas, suspensiones \\
    Periodo de consultas          & [Fecha inicio – Fecha acuerdo] \\
    Sindicatos firmantes          & [CCOO, UGT, FINE — verificar] \\
    Medidas sociales de acompañamiento & [Outplacement, complementos, formación] \\
    \bottomrule
  \end{tabularx}
\end{table}

\textbf{Valoración del proceso:}

[Descripción a completar — ¿fue percibido como justo? ¿Qué impacto tuvo
en el clima laboral de la plantilla que se quedó? ¿Survivor syndrome?]

\subsubsection{Downsizing y reestructuración continua}

% TODO: Más allá del ERE de 2021, ¿hay un proceso de reducción continua
% (no renovación de contratos, digitalización que sustituye puestos, etc.)?
% La automatización de procesos bancarios y la banca digital reducen la
% necesidad de ciertos perfiles de back-office.

[Descripción a completar]

\subsection{Recolocación: el \textit{outplacement}}

\subsubsection{Servicios de \textit{outplacement} en el ERE de CaixaBank}

% TODO: ¿Ofreció CaixaBank servicios de outplacement a los empleados
% que salieron? ¿Lo gestionaron internamente o con una consultora externa
% (p.ej. Lee Hecht Harrison, Randstad RiseSmart, Adecco)?
% ¿Cuántos empleados se beneficiaron? ¿Con qué tasa de recolocación?

El outplacement es un conjunto de servicios de apoyo a los empleados
que abandonan la empresa (voluntaria o involuntariamente), orientados a
facilitar su reincorporación al mercado laboral:

\begin{itemize}
  \item Orientación y coaching profesional.
  \item Actualización y elaboración del currículum.
  \item Preparación para entrevistas de trabajo.
  \item Identificación de oportunidades en el mercado.
  \item Apoyo psicológico durante la transición.
\end{itemize}

[Descripción de cómo se aplicó en CaixaBank]

\textbf{Ventajas del outplacement para CaixaBank:}
\begin{itemize}
  \item Mejora la imagen como empleador responsable (\textit{employer branding}).
  \item Reduce el conflicto y los litigios laborales.
  \item Mitiga el impacto emocional en los empleados que salen, lo que
        también reduce el síndrome del superviviente en los que se quedan.
\end{itemize}

\begin{notaequipo}
  Preguntad en la entrevista:
  \begin{enumerate}
    \item ¿Han recurrido a ERTE en algún momento (COVID, fusión)?
    \item ¿Qué condiciones se ofrecieron en el ERE de 2021?
          ¿Hubo jubilaciones anticipadas? ¿A partir de qué edad?
    \item ¿Se ofreció outplacement a los empleados afectados por el ERE?
          ¿Quién lo gestionó? ¿Cuáles fueron los resultados?
    \item ¿Cómo afectó el ERE al clima laboral de los empleados que se
          quedaron? ¿Tomaron medidas para gestionar ese impacto?
    \item ¿Realizan entrevistas de salida cuando un empleado dimite?
          ¿Qué hacen con esa información?
  \end{enumerate}
\end{notaequipo}
